% frontmatter/abstract-th.tex — บทคัดย่อภาษาไทย
\chapter*{บทคัดย่อ}
\addcontentsline{toc}{chapter}{บทคัดย่อภาษาไทย}

\noindent\textbf{หัวข้อปัญหาพิเศษ:}\ \thesistitleTH\\
\textbf{โดย:}\ \authorone, \authortwo, \authorthree\\
\textbf{อาจารย์ที่ปรึกษา:}\ \advisor\\
\textbf{ปีการศึกษา:}\ \academicyear

\vspace{0.6cm}

แพลตฟอร์มวิดีโอ YouTube มีผู้ใช้งานต่อเดือนมากกว่า 2.5 พันล้านคน และเป็นช่องทางการตลาด
ที่สำคัญของผู้สร้างเนื้อหาภาษาไทย การพยากรณ์ \emph{ความไวรัล} (virality) ของวิดีโอตั้งแต่
ก่อนเผยแพร่ จึงมีคุณค่าในเชิงพาณิชย์และเชิงวิชาการ งานวิจัยนี้กำหนดปัญหาเป็นการจำแนก
ประเภททวิภาคของวิดีโอภาษาไทยว่าจะเป็นไวรัลหรือไม่ โดยใช้ \emph{ดัชนีความไวรัลภายในช่อง}
(Per-channel Virality Index, $\mathit{VI}$) ซึ่งสร้างจากผลรวมแบบถ่วงน้ำหนักของอัตรา
ส่วนการมีส่วนร่วม (engagement rate) อัตราส่วนการกดถูกใจ และอัตราส่วนความคิดเห็น
หลังการแปลงด้วยลอการิทึมและ z-score ภายในช่องเจ้าของ ป้ายกำกับเชิงบวกถูกกำหนดเป็น
เดไซล์บนสุด (top-decile) ของแต่ละช่องเพื่อหลีกเลี่ยงไม่ให้ขนาดช่องเป็นปัจจัยรบกวน

ชุดข้อมูลประกอบด้วยวิดีโอจำนวน 23{,}431 รายการจาก 82 ช่องที่จัดเก็บผ่าน YouTube Data
API v3 มีอัตราส่วนคลาสบวก 10.16\% การแบ่งชุดข้อมูลใช้ \emph{แบบไวต่อเวลาผสมจัดกลุ่ม
ตามช่อง} (channel-grouped, time-aware split) โดยกำหนดให้ชุดทดสอบเป็นวิดีโอที่ตีพิมพ์
ในช่วง 15\% หลังสุดตามเวลา และทำการสุ่มลดข้อมูล (undersampling) ในอัตราส่วน 3:1
\emph{เฉพาะชุดฝึกฝนเท่านั้น} เพื่อรักษาการกระจายตัวตามธรรมชาติของชุดตรวจสอบและ
ชุดทดสอบ

งานวิจัยเปรียบเทียบแบบจำลองทรานส์ฟอร์เมอร์ภาษาไทยขนาดและสถาปัตยกรรมต่างกัน
3 แบบ ได้แก่ \textbf{WangchanBERTa} (105M, RoBERTa SentencePiece) \textbf{PhayaThaiBERT}
(278M, RoBERTa) และ \textbf{XLM-RoBERTa-large} (560M, multilingual) ทุกแบบจำลอง
ปรับจูนแบบเต็ม (full fine-tuning) บน GPU NVIDIA T4/A10G ผ่าน Hugging Face Jobs
แบบจำลองพื้นฐาน (Logistic Regression, LightGBM, XGBoost) ใช้ชุดคุณลักษณะเชิงโครงสร้าง
27 ตัว ผสมกับ TF-IDF อักขระ-N-gram และฟีเจอร์อารมณ์ (sentiment) จาก Wisesight-Sentiment
นอกจากนี้ยังพัฒนาแบบจำลอง \emph{ไฮบริด} (MLP + LightGBM head) ซึ่งใช้เวกเตอร์ฝัง
ข้อความรวมกับฟีเจอร์เชิงโครงสร้าง ปิดท้ายด้วย \emph{การประกอบแบบเรียงซ้อน} (stacking
ensemble) ที่ใช้ Logistic Regression เป็นตัวเรียนระดับเมตา (meta-learner) บนความน่าจะเป็น
ของฐาน 15 แบบจำลอง

จากการประเมินบนชุดทดสอบเดียวกัน แบบจำลองที่ดีที่สุดคือ \textbf{stacking ensemble
ที่ปรับเทียบด้วย Platt} ซึ่งให้ค่า ROC-AUC = \textbf{0.6914} โดยมีช่วงความเชื่อมั่นบูตสแตรป
95\% \textbf{[0.6625, 0.7174]}, ECE = 0.016 ซึ่งสูงกว่าแบบจำลองที่ดีที่สุดเดี่ยว
(LightGBM-Structured+TF-IDF, 0.6728) อย่างมีนัยสำคัญ ($p$ = $6.9\!\times\!10^{-53}$
จาก McNemar's test) ในระดับแบบจำลองทรานส์ฟอร์เมอร์เดี่ยว PhayaThaiBERT (0.6451)
ทำงานได้ดีที่สุด ตามด้วย WangchanBERTa (0.6402) และ XLM-RoBERTa-large (0.5450)
การทดสอบ Cochran's Q ระหว่างทรานส์ฟอร์เมอร์ทั้งสามให้ $Q = 612.69$,
$p = 9.0\!\times\!10^{-134}$ ($df = 2$) แสดงว่ามีความแตกต่างของรูปแบบความผิดพลาด
อย่างมีนัยสำคัญ การแบ่งย่อยรายช่องชี้ว่าแบบจำลองทำงานได้ดีกว่าในช่องขนาดใหญ่
(ROC-AUC = 0.7459 สำหรับ $\geq$1M ผู้ติดตาม) เทียบกับช่องขนาดกลาง (0.6273)

งานวิจัยรายงานผลลัพธ์ \emph{เชิงลบที่ตรงไปตรงมา} ว่าการเพิ่มข้อความหลายฟิลด์
(title + description + channel) ทำให้ประสิทธิภาพ PhayaThaiBERT \emph{ลดลง}
จาก 0.6451 เหลือ 0.5455 ซึ่งสะท้อนว่าคำอธิบายวิดีโอภาษาไทยส่วนใหญ่เป็นข้อความ
มาตรฐาน (boilerplate) การวิเคราะห์ความสามารถอธิบาย (explainability) ผ่าน SHAP,
LIME และ Attention Rollout แสดงว่าตัวสำคัญที่สุดในการพยากรณ์คือฟีเจอร์เชิงโครงสร้าง
(ความยาวชื่อ ความหนาแน่นของอีโมจิ ผู้ติดตามช่อง) ในขณะที่สัญญาณจากตัวภาษาเสริม
ผ่านการประกอบแบบเรียงซ้อน

\noindent\textbf{คำสำคัญ:} การพยากรณ์ความไวรัลของวิดีโอ, ทรานส์ฟอร์เมอร์ภาษาไทย,
WangchanBERTa, PhayaThaiBERT, XLM-RoBERTa, การประกอบแบบเรียงซ้อน,
McNemar's test, Cochran's Q, การปรับเทียบความน่าจะเป็น, SHAP, LIME

\newpage
